\section{Approach}\label{sec:approach}

The system follows the standard IR pipeline: preprocessing, indexing, query
processing, ranking, and evaluation. The implementation is built on the provided
Python template. Sentence segmentation uses NLTK Punkt, tokenization uses the
Penn Treebank tokenizer, inflection reduction uses Porter stemming, and
stopword removal uses the NLTK English stopword list. Because stemming is
performed before stopword removal, the implementation also removes stemmed
stopwords such as \emph{wa} and \emph{ha}. Tokens of length one and non-alphabetic
tokens are filtered to reduce noisy matches.

\subsection{Baseline TF-IDF Retrieval}

The baseline represents each document as a TF-IDF vector and ranks documents by
cosine similarity with the query. For a term $t$ in document $d$, the system uses
log-scaled term frequency:

\begin{equation}
TF(t,d)=1+\log(f_{t,d}),
\end{equation}

where $f_{t,d}$ is the raw count of $t$ in $d$. The implementation uses smoothed
IDF:

\begin{equation}
IDF(t)=\log\left(\frac{N+1}{df_t+1}\right)+1.
\end{equation}

The document weight is $TF(t,d)\times IDF(t)$. Query vectors use the same
weighting scheme, restricted to terms present in the document vocabulary. If a
query has no in-vocabulary terms, the system cannot compute a meaningful lexical
match and returns the default document order.

\subsection{Efficient Inverted TF-IDF Index}

Instead of comparing each query with every full document vector, the system
stores an inverted TF-IDF index from term to document weights. During ranking,
only documents sharing at least one query term receive a non-zero dot product.
The remaining documents are appended with score zero. This keeps the ranking
equivalent to cosine scoring for matching documents while avoiding unnecessary
work for documents that cannot match the query.

\subsection{Pseudo-Relevance Feedback}

The main retrieval improvement is Rocchio-style pseudo-relevance feedback. The
system first ranks documents using the original query. It then assumes that the
top retrieved document is relevant and expands the query vector using that
document's TF-IDF vector:

\begin{equation}
\vec{q'}=\alpha\vec{q}+\frac{\beta}{m}\sum_{d \in D_m}\vec{d},
\end{equation}

where $D_m$ is the set of top-$m$ feedback documents. In the implementation,
$m=1$, $\alpha=1.0$, and $\beta=0.2$. This setting keeps the original query
dominant while allowing related technical terms from the highest ranked document
to improve recall. The method is well suited to Cranfield because many relevant
documents describe the same aerodynamics concept using related but not identical
technical vocabulary.

\subsection{Limitations Addressed}

The improvements target three observed VSM limitations. First, stemming reduces
mismatches between related word forms, such as singular and plural forms or
technical suffix variants. Second, stemmed stopword removal prevents common
function words from surviving after stemming. Third, feedback expansion reduces
the vocabulary mismatch problem by adding terms from an initially retrieved
document. These changes do not solve true semantic ambiguity, but they improve
top-ranked retrieval while keeping the system interpretable.
